% Transcribed human main text; sources and scope in tex/008.
Suppose that we have a set of lines in $\mathbb R^3$. A joint of the set of lines is a point which lies in three non-coplanar lines. The joints problem asks what is the maximal number of joints that can be formed from $L$ lines.

\begin{theorem}[0.1]
Any $L$ lines in space determine $\leq 10L^{3/2}$ joints.
\end{theorem}
\begin{lemma}[Main Lemma]
If a set of lines has $J$ joints, then one of the lines contains $\leq 3J^{1/3}$ joints.
\end{lemma}
The main lemma implies the theorem by removing the lines one at a time. We start with $L$ lines and $J$ joints. By cutting out one line, we reduce the number of joints by $\leq 3J^{1/3}$. We look at the remaining lines $L-1$ lines, which contain $\leq J$ joints. One of the lines has $\leq 3J^{1/3}$ joints on it. Removing this line, we reduce the number of joints by $\leq 3J^{1/3}$. We remove all $L$ lines, one at a time. Each time the number of joints decreases by $\leq 3J^{1/3}$, and we end up with no joints. Therefore, $J\leq L(3J^{1/3})$. Rearranging we get $J^{2/3}\leq 3L$, which implies the theorem.

Now we turn to the proof of the main lemma.
\begin{proof}
Let $P$ be a lowest degree non-zero polynomial that vanishes at every joint. The degree of $P$ is $\leq 3J^{1/3}$ by dimension counting, as in Lecture 1. If every line has $>3J^{1/3}$ joints, then $P$ must vanish on every line.

\begin{lemma}[0.2]
If $x$ is a joint lying in three (non-coplanar) lines, and if a smooth function $F:\mathbb R^3\to\mathbb R$ vanishes on the lines, then $\nabla F$ vanishes at $x$.
\end{lemma}
\begin{proof}
Let $v_1,v_2,v_3$ be tangent vectors for the three lines. The directional derivative of $F$ in the direction $v_i$ must vanish at $x$. So we have $\nabla F(x)\cdot v_i=0$ for each $i$. Since the $v_i$ are a basis of $\mathbb R^3$, we have $\nabla F(x)=0$.
\end{proof}
So we see that the derivates of $P$ vanish at each joint. The derivates have smaller degree than $P$. Since $P$ was a minimal degree non-zero polynomial that vanishes at each joint, the derivatives of $P$ are all the zero polynomial! Then $P$ must be constant, and we get a contradiction.
\end{proof}

\paragraph{所引用的维数计数结果（Lecture 1，pp.~1--2）}
\begin{corollary}[0.2]
If $\dim V(d)>s$, then there is a non-zero polynomial $P$ of degree $\leq d$ that vanishes on the given finite set.
\end{corollary}
The dimension of $V(d)$ is $\binom{d+n}{n}\geq d^n/n!$. For $n$ fixed and $d$ large, $d^n/n!$ is a good approximation of the dimension. Therefore, we get the following corollary.
\begin{corollary}[0.3]
For any set of $s$ points in $F^n$, there is a non-zero polynomial that vanishes on the set with degree $\leq ns^{1/n}$.
\end{corollary}
